% Учебный черновик по PySATL Core #128.
% Компиляция: pdfLaTeX с пакетами кириллицы. Автор и руководитель уточняются перед сдачей.
\documentclass[12pt,a4paper]{article}
\usepackage[T2A]{fontenc}
\usepackage[utf8]{inputenc}
\usepackage[russian,english]{babel}
\usepackage[margin=25mm]{geometry}
\usepackage{microtype}
\usepackage{hyperref}
\usepackage{xurl}
\usepackage{array}
\usepackage{longtable}
\usepackage{booktabs}
\usepackage{enumitem}
\hypersetup{colorlinks=true,linkcolor=blue,urlcolor=blue,citecolor=blue}
\setlength{\parindent}{1.25cm}
\setlength{\parskip}{0.35em}
\sloppy

\title{Архитектура вычислений характеристик и численных аппроксимаций\\ преобразованных распределений в PySATL Core}
\author{Исполнитель: [ФИО уточняется]\\ Руководитель: [ФИО уточняется]}
\date{Черновик для согласования, 25 октября 2026 г.}

\begin{document}
\maketitle
\begin{center}
\textbf{Основание:} \url{https://github.com/PySATL/pysatl-core/issues/128}\\
\textbf{Кодовый репозиторий:} \url{https://github.com/PySATL/pysatl-core}\\
\textbf{Репозиторий отчёта:} \url{https://github.com/siwreienta/pysatl-core-report}.
\end{center}

\section*{Аннотация}
В проекте PySATL Core уже реализованы параметрические семейства распределений, граф вычисления характеристик, аффинные и бинарные преобразования, конечные смеси и первые средства численной интерполяции PDF, CDF и PPF. Однако методы, которые распределение предоставляет как начальные точки вычисления, сейчас называются \texttt{analytical\_computations}, хотя в эту карту попадают также вычисления преобразованных и аппроксимированных распределений. Кроме того, методы родителей в преобразованиях разрешаются при создании дочернего объекта. Это затрудняет ясное определение момента выбора метода, его подготовки, повторного использования и оценки численной ошибки.

Цель работы — разработать согласованные контракты входных вычислений характеристик и численных аппроксимаций для PySATL Core, а затем проверить их на одном ограниченном методе арифметики независимых одномерных непрерывных случайных величин. Сначала будут исследованы существующие классы и стратегия вычислений, определена семантика \texttt{entry\_computations} и принято архитектурное решение о признаке \texttt{is\_analytical}. Затем планируется выделить самостоятельный модуль аппроксимаций, описать жизненный цикл материализованного распределения, подготовку и кэширование методов, область аппроксимации и диагностические метаданные. Существующие интерполяционные прототипы будут перенесены с документированным контрактом ошибок для PDF, CDF и PPF. Отдельно будет проверено, нужен ли поздний выбор методов родителей при построении цепочки преобразований.

В численной части предполагается выбрать одну операцию и конечный набор тестовых распределений, построить сегментное представление результата и сравнить его с доступными аналитическими эталонами. Проверка включает вероятностные инварианты, ошибки на узлах и между ними, поведение у границ и возможность повторно использовать полученное представление. Работа опирается на идеи PaCAL и публикации о численной арифметике распределений, но не предполагает копирование их API или реализации. Ожидаемый результат — обоснованная архитектура, воспроизводимый эксперимент, тесты и документация с явно указанными ограничениями. Предлагаемое название и окончательный объём подлежат согласованию с преподавателем.

\section{Постановка и границы работы}
Исходный epic \cite{epic} ставит задачу перехода от термина \texttt{analytical\_computations} к \texttt{entry\_computations}, выделения модуля аппроксимаций, пересмотра жизненного цикла преобразований и добавления первого метода по мотивам PaCAL. Достижимая цель учебной работы: спроектировать эти контракты и продемонстрировать их на одном методе для независимых одномерных непрерывных распределений, сохранив существующие аффинные, бинарные и смесевые сценарии либо описав необходимую миграцию.

Критерии результата: принятое ADR; согласованное API и тесты; отдельный модуль аппроксимаций с описанными fit/cache, доменом и диагностикой; проверка выбранного метода на эталонных случаях. Точные допуски численной ошибки должны быть определены после базового эксперимента и зафиксированы до окончательной валидации. Не рассматриваются в первом срезе зависимые, многомерные и смешанные распределения, все виды сингулярностей, универсальный benchmark-фреймворк \cite{benchissue} и рефакторинг численного runtime \cite{runtimeissue}.

\section{Задачи и ориентировочные этапы}
Обзор литературы формирует постановку и обоснование, но не заменяет задачи реализации. Если работа займёт весь учебный год, этапы могут распределяться так:
\begin{enumerate}[leftmargin=*]
  \item \textbf{Осень--зима 2026:} определить семантику \texttt{entry\_computations} и \texttt{is\_analytical}, затем выполнить миграцию API, тестов, примеров и документации вместе с ADR в одном PR.
  \item \textbf{Зима 2026--27:} выделить самостоятельный модуль аппроксимаций, перенести PDF/CDF/PPF-прототипы и закрепить контракты материализации, подготовки, кэширования и ошибок в одном PR.
  \item \textbf{Зима--весна 2027:} сравнить ранний и поздний выбор родительских методов, уточнить роль \texttt{LightweightDistribution} и выполнить обоснованный рефакторинг преобразований с тестами в отдельном PR.
  \item \textbf{Весна 2027:} реализовать сегментное представление и один ограниченный численный метод, проверить эталонные случаи и вероятностные инварианты, добавить воспроизводимый пример и документацию в одном PR.
\end{enumerate}
Каждая техническая задача объединяет необходимое исследование с реализацией и завершается проверяемым PR в PySATL Core. Сроки являются ориентирами для обсуждения.

\section{Обзор источников и аналогов}
\subsection{Математические подходы}
Jaroszewicz и Korzeń \cite{jaroszewicz2012} исследуют численную арифметику независимых случайных величин и эффективную аппроксимацию плотностей полиномами Чебышёва. Это прямое основание для выбора независимого одномерного случая, но их результаты нельзя автоматически распространять на зависимые величины или на любую сингулярность. PaCAL \cite{pacal} демонстрирует программную реализацию похожих идей с пригодностью результата для последующих операций. Для PySATL его полезно рассматривать как архитектурный аналог, а не зависимость: PySATL имеет собственный граф характеристик, стратегии и иные контракты.

Книга Drew и соавторов \cite{drew2017} показывает более широкий класс вычислительных операций над распределениями, включая подход APPL. Поэтому новизна данной работы должна формулироваться как интеграция проверяемых численных контрактов в PySATL Core, а не как изобретение операций над распределениями. Книга Trefethen \cite{trefethen2019} даёт теоретические ориентиры для интерполяции Чебышёва и выбора сегментов, а книга Higham \cite{higham2002} — язык для анализа округления и устойчивости. На момент черновика книги изучены по официальным описаниям и доступным разделам; полный текст необходимо прочитать для окончательной версии обзора.

\subsection{Существующие инструменты и реализация PySATL}
В SciPy интерфейс \texttt{rv\_continuous} допускает определение части характеристик с численным выводом остальных \cite{scipyrv}. Это полезный ориентир для сравнения API, однако PySATL требуется явный граф зависимостей и выбор между методами с метками. \texttt{PchipInterpolator} сохраняет монотонность и уже используется в прототипах CDF/PPF \cite{scipypchip}. Он решает локальную задачу интерполяции, но не задаёт согласованность независимо приближённых PDF и CDF. \texttt{scipy.integrate.quad} принимает точки возможных трудностей и возвращает оценку ошибки интеграла \cite{scipyquad}; \texttt{Chebyshev.interpolate} даёт готовый примитив на конечном интервале \cite{numpycheb}. Ни один из этих инструментов сам по себе не определяет, где в PySATL хранить breakpoints, как передавать их через преобразования и как сообщать пользователю итоговую ошибку.

Текущий код PySATL \cite{core} имеет \texttt{Distribution}, \texttt{DerivedDistribution}, стратегию, граф и первые интерполяции. PDF-прототип использует равномерную сетку на конечном интервале и обрезает отрицательные значения; контракт результата пока не содержит диагностики массы. Это делает его полезной базой сравнения, но требует проверки нормировки и поведения вне сетки.

\section{Анализ обзора и выводы для проекта}
\begin{longtable}{p{0.20\textwidth}p{0.34\textwidth}p{0.36\textwidth}}
\toprule
Подход & Сильная сторона для задачи & Ограничение и решение для PySATL \\
\midrule
\endfirsthead
\toprule
Подход & Сильная сторона для задачи & Ограничение и решение для PySATL \\
\midrule
\endhead
PaCAL и статья \cite{jaroszewicz2012,pacal} & Численная арифметика и сегментные представления, которые можно применять повторно & Переносить принципы и проверяемые случаи; разработать Core-native контракт, не копировать классы или код. \\
APPL \cite{drew2017} & Широкая алгоритмическая постановка операций над распределениями & Ограничить реализацию первым численным случаем и корректно описать новизну. \\
SciPy и NumPy \cite{scipyrv,scipypchip,scipyquad,numpycheb} & Готовые интерполяторы, интегратор и базовые probabilistic API & Добавить собственные договорённости о метаданных, домене, fit/cache и ошибке всей операции. \\
Текущий PySATL \cite{core} & Существующие стратегии, граф, трансформации и прототипы & Сохранить инварианты и совместимость; устранить терминологическую и модульную неоднозначность. \\
\bottomrule
\end{longtable}

Из сравнения следуют три решения для постановки. Во-первых, сначала следует установить семантику точки входа и аналитичности, иначе новый метод будет встроен в неоднозначный контракт. Во-вторых, качество PDF, CDF и PPF нужно задавать раздельно: независимая аппроксимация каждой характеристики допустима, но их взаимная согласованность не возникает автоматически. В-третьих, первый метод должен иметь явные ограничения по независимости, размерности и области определения. Этим уточнена исходная широкая цель epic: учебный результат не обещает полной реализации PaCAL, зато содержит законченный проверяемый срез.

\section{Начало первой задачи: аудит контрактов}
Локально изучены \texttt{Distribution}, \texttt{DerivedDistribution}, \texttt{ApproximatedDistribution}, \texttt{TransformationMethod}, \texttt{LightweightDistribution}, граф, \texttt{DefaultComputationStrategy} и исходные PDF/CDF/PPF-прототипы. Выявлено следующее:
\begin{itemize}[leftmargin=*]
  \item карта \texttt{analytical\_computations} обязательна и непуста, но у производного распределения не обязательно означает полностью аналитическую цепочку;
  \item \texttt{TransformationMethod} разрешает родительские методы при создании, а аналитичность родителя проверяет по первому варианту метки — это требует отдельного теста при нескольких вариантах;
  \item граф различает аналитические и трансформационные петли, но текущая стратегия выбирает первую доступную петлю без прямой фильтрации по \texttt{is\_analytical};
  \item \texttt{ApproximatedDistribution} хранит материализованные методы как самостоятельный объект с \texttt{bases=\{\}} и неаналитическими петлями.
\end{itemize}
Артефакт первой задачи — черновик \texttt{architecture/ADR-0001-entry-computations.md} с контекстом и предлагаемым решением. Для завершения задачи нужны таблица обращений к старому имени, тест на несколько меток, согласованное ADR и миграция API, тестов и документации в одном PR. Сравнение eager/lazy с проверкой числа fit-вызовов войдёт в отдельную задачу о жизненном цикле преобразований. Сложность выше простого переименования: имя связано с графом, стратегией, снимками родителей и кешированием. Кодовый репозиторий пока не изменён.

\section*{Статус сдачи}
Название, комментарии к источникам и точные численные допуски остаются предметом согласования и последующей проверки. Отчёт размещается в отдельном GitHub-репозитории; дальнейшие этапы будут сдаваться ссылками на PR.

\begin{thebibliography}{99}
\bibitem{epic} PySATL Core. Epic \#128: Refactor entry computations, transformations, and approximations. \url{https://github.com/PySATL/pysatl-core/issues/128}.
\bibitem{core} PySATL Core, исходный код и документация. \url{https://github.com/PySATL/pysatl-core}.
\bibitem{jaroszewicz2012} Jaroszewicz S., Korzeń M. Arithmetic Operations on Independent Random Variables: A Numerical Approach. SIAM Journal on Scientific Computing, 34(3), A1241--A1265, 2012. \url{https://doi.org/10.1137/110839680}.
\bibitem{pacal} PaCAL: ProbAbilistic CALculator. Авторский репозиторий. \url{https://github.com/jszymon/pacal}.
\bibitem{drew2017} Drew J. H., Evans D. L., Glen A. G., Leemis L. M. Computational Probability, 2nd ed. Springer, 2017. \url{https://link.springer.com/book/10.1007/978-3-319-43323-3}.
\bibitem{trefethen2019} Trefethen L. N. Approximation Theory and Approximation Practice, Extended Edition. SIAM, 2019. \url{https://epubs.siam.org/doi/book/10.1137/1.9781611975949}.
\bibitem{higham2002} Higham N. J. Accuracy and Stability of Numerical Algorithms, 2nd ed. SIAM, 2002. \url{https://epubs.siam.org/doi/book/10.1137/1.9780898718027}.
\bibitem{scipyrv} SciPy. \texttt{rv\_continuous}. \url{https://docs.scipy.org/doc/scipy/reference/generated/scipy.stats.rv_continuous.html}.
\bibitem{scipypchip} SciPy. \texttt{PchipInterpolator}. \url{https://docs.scipy.org/doc/scipy/reference/generated/scipy.interpolate.PchipInterpolator.html}.
\bibitem{scipyquad} SciPy. \texttt{integrate.quad}. \url{https://docs.scipy.org/doc/scipy/reference/generated/scipy.integrate.quad.html}.
\bibitem{numpycheb} NumPy. \texttt{Chebyshev.interpolate}, версия документации 2.1. \url{https://numpy.org/doc/2.1/reference/generated/numpy.polynomial.chebyshev.Chebyshev.interpolate.html}.
\bibitem{runtimeissue} PySATL Core. Issue \#127: Configurable numerical runtime and array backends. \url{https://github.com/PySATL/pysatl-core/issues/127}.
\bibitem{benchissue} PySATL Core. Issue \#124: Benchmarking framework for performance and numerical quality. \url{https://github.com/PySATL/pysatl-core/issues/124}.
\end{thebibliography}
\end{document}
